\chapter{晶格、电子气与能带的核心推导}

\section{倒格基矢与体积}
要求
$\vb a_i\cdot\vb b_j=2\pi\delta_{ij}$。$\vb b_1$ 必须垂直
$\vb a_2,\vb a_3$，故
$\vb b_1=C\vb a_2\times\vb a_3$。点乘 $\vb a_1$：
\[
C\,\vb a_1\cdot(\vb a_2\times\vb a_3)=2\pi,
\]
\[
\boxed{\vb b_1=2\pi
\frac{\vb a_2\times\vb a_3}{\Omega_c}},
\quad
\Omega_c=\vb a_1\cdot(\vb a_2\times\vb a_3).
\]
循环置换得其余基矢。矩阵形式 $B^TA=2\pi I$：
\[
\boxed{\Omega_c^*=\det B=\frac{(2\pi)^3}{\Omega_c}}.
\]

\section{Laue 与 Bragg 等价}
格点散射振幅
\[
\mathcal A\propto\sum_{\vb R}
e^{i(\vb k-\vb k')\cdot\vb R}.
\]
所有格点同相要求
\[
(\vb k'-\vb k)\cdot\vb R=2\pi n
\quad\forall\vb R,
\]
所以
\[
\boxed{\vb k'-\vb k=\vb G}.
\]
弹性散射 $|\vb k|=|\vb k'|=2\pi/\lambda$。几何上
$|\vb G|=2k\sin\theta$；晶面族
$|\vb G|=2\pi n/d$：
\[
2\frac{2\pi}{\lambda}\sin\theta=\frac{2\pi n}{d}
\Rightarrow\boxed{2d\sin\theta=n\lambda}.
\]

\section{单原子链色散}
\[
M\ddot u_n=K(u_{n+1}+u_{n-1}-2u_n).
\]
取
$u_n=u e^{i(kna-\omega t)}$：
\[
-M\omega^2u
=Ku(e^{ika}+e^{-ika}-2)
=-4Ku\sin^2(ka/2).
\]
所以
\[
\boxed{\omega=2\sqrt{K/M}|\sin(ka/2)|}.
\]
$ka\ll1$：
\[
\omega\simeq a\sqrt{K/M}|k|,\quad
v_s=a\sqrt{K/M}.
\]
$k=0$ 是刚性平移，故频率必须为零。

\section{Debye 态密度与热容}
三维 $q$ 空间每态体积 $(2\pi)^3/V$。含三偏振：
\[
N(q)=3\frac{V}{(2\pi)^3}\frac{4\pi q^3}{3}
=\frac{Vq^3}{2\pi^2}.
\]
$\omega=v_sq$：
\[
g(\omega)=\dv{N}{\omega}
=\frac{3V\omega^2}{2\pi^2v_s^3}.
\]
以 $\int_0^{\omega_D}g\,\dd\omega=3N$ 定义截止。热能
\[
U=9N\kb T\left(\frac T{\Theta_D}\right)^3
\int_0^{\Theta_D/T}\frac{x^3}{e^x-1}\dd x.
\]
$T\ll\Theta_D$ 时上限延至无穷，
$\int_0^\infty x^3/(e^x-1)\dd x=\pi^4/15$：
\[
U=\frac{3\pi^4}{5}N\kb\frac{T^4}{\Theta_D^3},
\quad
\boxed{C_V=\frac{12\pi^4}{5}N\kb(T/\Theta_D)^3}.
\]

\section{自由电子态密度与费米能}
含自旋的 $k$ 球内态数
\[
N(k)=2\frac V{(2\pi)^3}\frac{4\pi k^3}{3}
=\frac{Vk^3}{3\pi^2}.
\]
$E=\hbar^2k^2/2m$：
\[
\dv kE=\frac m{\hbar^2k},
\]
\[
g(E)=\dv Nk\dv kE
=\frac{Vk^2}{\pi^2}\frac m{\hbar^2k}
=\boxed{\frac V{2\pi^2}
\left(\frac{2m}{\hbar^2}\right)^{3/2}\sqrt E}.
\]
$T=0$：
\[
N=\int_0^{E_F}g(E)\dd E
=\frac V{3\pi^2}
\left(\frac{2mE_F}{\hbar^2}\right)^{3/2},
\]
\[
\boxed{E_F=\frac{\hbar^2}{2m}(3\pi^2n)^{2/3}}.
\]

\section{Bloch 定理}
周期势使
$[H,T_{\vb R}]=0$。平移算符幺正，取共同本征态：
\[
T_{\vb R}\psi=\lambda(\vb R)\psi,\quad|\lambda|=1.
\]
平移群性质
$T_{\vb R}T_{\vb R'}=T_{\vb R+\vb R'}$ 给
\[
\lambda(\vb R)\lambda(\vb R')
=\lambda(\vb R+\vb R').
\]
其一维表示写成
$\lambda=e^{i\vb k\cdot\vb R}$。定义
\[
u_{\vb k}(\vb r)=e^{-i\vb k\cdot\vb r}\psi_{\vb k}(\vb r),
\]
\[
u(\vb r+\vb R)
=e^{-i\vb k\cdot(\vb r+\vb R)}
e^{i\vb k\cdot\vb R}\psi(\vb r)=u(\vb r).
\]
故
\[
\boxed{\psi_{n\vb k}=e^{i\vb k\cdot\vb r}u_{n\vb k}}.
\]

\section{区边界能隙}
周期势
$V=\sum_GV_Ge^{iGx}$ 耦合 $|k\rangle$ 和 $|k-G\rangle$。
$k=G/2$ 时二者自由能相同。简并子空间：
\[
\begin{pmatrix}
E_0+V_0&V_G\\
V_G^*&E_0+V_0
\end{pmatrix}.
\]
特征方程
\[
(E_0+V_0-E)^2-|V_G|^2=0,
\]
\[
E_\pm=E_0+V_0\pm|V_G|,
\quad\boxed{\Delta=2|V_G|}.
\]
本征态是两平面波的相干叠加，即不同相位的驻波；其电荷密度与离子势重叠不同，故能量劈裂。

\section{Boltzmann 电导}
均匀稳态、弱电场：
\[
\dot{\vb k}\cdot\nabla_kf
=-\frac{f-f_0}{\tau},\qquad
\hbar\dot{\vb k}=-e\vb E.
\]
令 $f=f_0+\delta f$ 并只保留一阶：
\[
-\frac e\hbar\vb E\cdot\nabla_kf_0
=-\frac{\delta f}{\tau}.
\]
因
$\nabla_kf_0=(\partial_Ef_0)\nabla_kE
=\hbar\vb v\partial_Ef_0$：
\[
\delta f=e\tau\vb E\cdot\vb v(-\partial_Ef_0).
\]
\[
j_\alpha=-e\int\frac{\dd^3k}{(2\pi)^3}
v_\alpha\delta f
=\sigma_{\alpha\beta}E_\beta,
\]
\[
\boxed{\sigma_{\alpha\beta}
=e^2\int\frac{\dd^3k}{(2\pi)^3}
\tau v_\alpha v_\beta(-\partial_Ef_0)}.
\]
零温 $-\partial_Ef_0=\delta(E-E_F)$，明确显示输运是费米面性质。

\section{半导体载流子浓度}
导带底抛物线态密度
\[
g_c(E)=\frac1{2\pi^2}
\left(\frac{2m_e^*}{\hbar^2}\right)^{3/2}
\sqrt{E-E_c}.
\]
非简并极限
$f\simeq e^{-(E-E_F)/\kb T}$：
\[
n=\int_{E_c}^\infty g_c(E)f(E)\dd E.
\]
令 $\epsilon=E-E_c$：
\[
n=e^{-(E_c-E_F)/\kb T}
\frac1{2\pi^2}\left(\frac{2m_e^*}{\hbar^2}\right)^{3/2}
\int_0^\infty\epsilon^{1/2}e^{-\epsilon/\kb T}\dd\epsilon.
\]
用
$\int_0^\infty x^{1/2}e^{-x}\dd x=\sqrt\pi/2$：
\[
\boxed{n=N_ce^{-(E_c-E_F)/\kb T}},
\quad
N_c=2\left(\frac{m_e^*\kb T}{2\pi\hbar^2}\right)^{3/2}.
\]
同理
$p=N_ve^{-(E_F-E_v)/\kb T}$，相乘：
\[
\boxed{np=N_cN_ve^{-E_g/\kb T}=n_i^2}.
\]

\section{Curie--Weiss 平均场}
独立两能级磁矩
\[
Z=2\cosh(\beta\mu B),\quad
\langle m\rangle=\frac1\beta\pdv{\ln Z}{B}
=\mu\tanh(\beta\mu B).
\]
弱场
$\langle m\rangle\simeq\mu^2B/\kb T$，得 Curie 定律。
平均场取
$H_{\rm eff}=H+\lambda M$：
\[
M=\frac CT(H+\lambda M).
\]
移项：
\[
M\left(1-\frac{C\lambda}{T}\right)=\frac CT H,
\]
\[
\boxed{\chi=\frac MH=\frac C{T-C\lambda}
=\frac C{T-T_c}}.
\]

\section{London 穿透深度}
London 第二方程
\[
\nabla\times\vb j_s=-\frac{n_se^{*2}}{m^*}\vb B.
\]
对 Ampère 方程
$\nabla\times\vb B=\mu_0\vb j_s$ 取旋度：
\[
\nabla\times\nabla\times\vb B
=\mu_0\nabla\times\vb j_s.
\]
用
$\nabla\times\nabla\times\vb B
=\nabla(\nabla\cdot\vb B)-\nabla^2\vb B=-\nabla^2\vb B$：
\[
-\nabla^2\vb B
=-\mu_0\frac{n_se^{*2}}{m^*}\vb B.
\]
\[
\boxed{\nabla^2\vb B=\vb B/\lambda_L^2},
\quad
\lambda_L=\sqrt{\frac{m^*}{\mu_0n_se^{*2}}}.
\]
半无限样品 $x>0$ 的有限解
$B(x)=B(0)e^{-x/\lambda_L}$。

\begin{selftestbox}
\begin{enumerate}
\item 倒格原胞体积如何由矩阵行列式推出？
\item Laue 条件怎样化成 Bragg 条件？
\item 单原子链周期色散来自哪个离散结构？
\item Debye 态密度的 $\omega^2$ 来自哪里？
\item 自由电子态密度如何包含自旋因子？
\item Bloch 平移本征值为何是指数形式？
\item 能隙为何等于 $2|V_G|$？
\item Boltzmann 电导中费米面窗口是什么？
\item 有效态密度 $N_c$ 的 $T^{3/2}$ 怎样出现？
\item London 方程如何区别超导与完美导体？
\end{enumerate}
\end{selftestbox}
